% ============================================================================
% Fichier  : partie1/P01_C02_preuve_chaine_possession.tex
% Titre    : La Preuve Numérique et la Chaîne de Possession
% Auteur   : MINKA MI NGUIDJOÏ Thierry Emmanuel
% Version  : 1.0 -- novembre 2024
% Corrigés : corriges/Corr_C02.tex
% ============================================================================
% Sources : ISO/IEC 27037:2012 (paraphrasée) ; loi n° 2010/012, art. 53 et 54
% (texte officiel, cité) ; FIPS 180-4 et état de l'art sur les collisions
% MD5/SHA-1 (présenté sans citation d'article).
% Faits nouveaux (bible) : chaîne de possession de la capture mémoire
% (scellé NC-01) ; trou de 17 h 40 dans la nuit du 14 au 15 mars ; scellés
% NC-01 à NC-06.
% ============================================================================

\chapter{La Preuve Numérique et la Chaîne de Possession}
\label{chap:preuve}
\chaptermark{Preuve et chaîne de possession}

\epigraph{%
    « Une pièce dont on ne peut pas dire où elle était\\
    ne vaut pas mieux qu'une pièce perdue. »%
}{Principe du cours}

\niveauChapitre{Fondamental}{%
    Maîtrise du Chapitre~\ref{chap:investigation}
}

\begin{boxobjectifs}
\begin{itemize}
    \item Distinguer la recevabilité d'une preuve et sa force probante.
    \item Expliquer ce qu'une empreinte cryptographique prouve, et ce
          qu'elle ne prouve pas.
    \item Choisir une fonction de hachage adaptée et justifier ce choix.
    \item Tenir une chaîne de possession complète, du recueil à la
          restitution.
    \item Identifier une rupture de la chaîne et en mesurer les
          conséquences.
    \item Appliquer les articles~53 et~54 de la loi n°~2010/012 à une
          perquisition portant sur des données.
\end{itemize}
\end{boxobjectifs}

\bigskip

% ============================================================================
\section*{Avant de commencer : un disque sur un bureau}
% ============================================================================

Bafoussam, samedi 15 mars 2025, 9~h~30. Mme~MBARGA demande où se trouve le
disque externe qui contient la capture mémoire du contrôleur de domaine. Il
est sur le bureau de M.~KAMGA, dans une enveloppe kraft fermée par un ruban
adhésif. Personne ne l'a mis sous clé pendant la nuit. La porte du bureau
du DSI ne ferme pas, et l'équipe de ménage passe le vendredi soir.

L'empreinte de la capture avait été calculée la veille à 10~h~41 et inscrite
au registre. Mme~MBARGA la recalcule devant M.~KAMGA et Mme~NDOUMBE~: elle
est identique. Le contenu n'a donc pas changé. Mais l'avocat de la défense,
dans dix-huit mois, posera une autre question~: pendant ces dix-sept heures,
qui pouvait toucher ce disque~?

\begin{boxalert}
\textbf{Violation typique.} Croire qu'une empreinte identique suffit à
défendre une pièce, et négliger la trace écrite de ceux qui l'ont détenue.

\smallskip
\textbf{Règle~:} l'empreinte prouve que le contenu n'a pas changé entre deux
calculs~; la chaîne de possession prouve qui détenait la pièce, où, et dans
quelles conditions. Les deux se complètent et aucune ne remplace l'autre.
\end{boxalert}

\separateur

% ============================================================================
\section{Recevable et Probante}
\label{sec:c02-recevable}
% ============================================================================

Deux questions se posent devant toute pièce, et elles ne se confondent pas.

\begin{boxdef}
\textbf{\termecle{Recevabilité}}~: aptitude d'une pièce à être versée au
débat. Une pièce obtenue en violation de la loi, par exemple par un accès
non autorisé à un système, peut être écartée.

\smallskip
\textbf{\termecle{Force probante}}~: poids que le décideur accorde à une
pièce recevable. Elle dépend de la fiabilité de la source, de l'intégrité de
la pièce, de la qualité de la méthode et de la clarté de son interprétation.
\end{boxdef}

En matière pénale, le droit camerounais, comme la plupart des droits issus
de la tradition civiliste, admet en principe la preuve par tout moyen et
laisse le juge apprécier librement la valeur des pièces [disposition du Code
de procédure pénale à vérifier par le lecteur sur le texte officiel]. Ce
principe ouvre la porte à la preuve numérique, mais il ne la rend pas
automatiquement convaincante. Une pièce numérique contestée sur son
intégrité ou sur sa provenance sera recevable et peu probante~: le juge la
lira, puis lui accordera peu de poids. L'enquêteur travaille donc moins pour
la recevabilité, que la loi accorde largement, que pour la force probante,
qu'il doit construire.

\begin{tableaucours}{L{4.0cm}L{8.2cm}}{%
    \tableauHeader{Ce qui fonde la force probante} &
    \tableauHeader{Ce que l'enquêteur doit pouvoir montrer}
}
Provenance & De quel système, de quel emplacement, à quelle heure la pièce
a été tirée \\
Intégrité & Que le contenu examiné est celui qui a été recueilli \\
Continuité & Qui a détenu la pièce, à chaque instant, et dans quelles
conditions \\
Méthode & Comment la pièce a été acquise et examinée, avec quels outils et
quelles versions \\
Interprétation & Ce que la pièce signifie, et les autres explications
possibles \\
\end{tableaucours}
\vspace{-0.8em}
\begin{center}{\small\itshape Tableau~2.1. Les fondements de la force probante}\end{center}

% ============================================================================
\section{L'Empreinte Cryptographique}
\label{sec:c02-empreinte}
% ============================================================================

\begin{boxdef}
\textbf{\termecle{Fonction de hachage cryptographique}}~: fonction qui
associe à une donnée de taille quelconque une valeur de taille fixe,
l'\textbf{empreinte} (\textit{hash}), de sorte qu'il soit pratiquement
impossible de trouver deux données différentes de même empreinte
(résistance aux collisions), ou de construire une donnée qui ait une
empreinte donnée (résistance aux préimages).
\end{boxdef}

Si l'on modifie un seul bit d'une image disque de 500~Go, son empreinte
SHA-256 change entièrement. Recalculer l'empreinte à chaque étape et
constater qu'elle n'a pas varié démontre donc que le contenu est resté le
même, sans avoir à comparer les deux copies octet par octet.

\subsection{Quelle fonction choisir}

\begin{tableaucours}{L{2.0cm}cL{4.6cm}L{3.8cm}}{%
    \tableauHeader{Fonction} &
    \tableauHeader{Taille} &
    \tableauHeader{État} &
    \tableauHeader{Usage en enquête}
}
MD5 & 128 bits & Collisions calculables depuis 2004~; on fabrique deux
fichiers distincts de même empreinte & Complément pour comparer à
d'anciennes bases, jamais seule \\
SHA-1 & 160 bits & Collision pratique publiée en 2017 & Complément
seulement \\
SHA-256 & 256 bits & Aucune attaque pratique connue & Empreinte de
référence \\
SHA-512, SHA-3 & 512 bits et plus & Aucune attaque pratique connue &
Possibles, moins répandues dans les outils \\
\end{tableaucours}
\vspace{-0.8em}
\begin{center}{\small\itshape Tableau~2.2. Fonctions de hachage et usage en enquête}\end{center}

La faiblesse de MD5 porte sur les collisions, pas directement sur les
préimages~: un attaquant peut fabriquer deux fichiers de même empreinte, il
ne peut pas, à ce jour, fabriquer un fichier qui ait l'empreinte d'un
fichier donné. Une image acquise avec sa seule empreinte MD5 reste donc
défendable en pratique, mais elle offre à la partie adverse un argument
qu'une empreinte SHA-256 lui refuse. On calcule aujourd'hui SHA-256, et l'on
ajoute MD5 quand un outil ou une base de référence l'exige.

\subsection{Ce que l'empreinte ne prouve pas}

\begin{itemize}
    \item Elle ne prouve pas que la copie est complète~: une acquisition qui
          a sauté des secteurs illisibles a une empreinte stable et un
          contenu lacunaire. Le journal de l'outil d'acquisition le dit.
    \item Elle ne prouve pas que la pièce vient du système annoncé~: c'est
          le rôle de la provenance, consignée au recueil.
    \item Elle ne prouve rien sur la période qui précède son premier
          calcul~: une capture modifiée avant d'être hachée garde, ensuite,
          une empreinte parfaitement stable.
\end{itemize}

La troisième limite explique pourquoi on calcule l'empreinte le plus tôt
possible, idéalement pendant l'acquisition, et pourquoi on la fait
constater par une seconde personne.

% ============================================================================
\section{La Chaîne de Possession}
\label{sec:c02-chaine}
% ============================================================================

\begin{boxdef}
\textbf{\termecle{Chaîne de possession}} (\textit{chain of custody})~:
documentation continue qui retrace, pour chaque pièce, chaque personne qui
l'a détenue, chaque transfert, chaque lieu de conservation et chaque
opération réalisée, depuis le recueil jusqu'à la restitution ou la
destruction.
\end{boxdef}

L'ISO/IEC~27037~\cite{iso27037_2012} demande que cette documentation permette
de répondre, pour chaque pièce, à des questions simples~: qui l'a recueillie,
quand, où, comment et pourquoi~; qui l'a détenue ensuite et à quel moment
elle a changé de mains~; quelles opérations ont été faites sur elle. On
paraphrase ici ces exigences sous la forme d'un formulaire.

\begin{boxprocedure}
\textbf{Le formulaire de chaîne de possession d'une pièce}
\begin{enumerate}
    \item Identifiant unique de la pièce et numéro de scellé.
    \item Description~: nature, marque, modèle, numéro de série, capacité.
    \item Provenance~: système, emplacement physique, propriétaire.
    \item Recueil~: date, heure, lieu, personne, méthode, outil et version.
    \item Empreintes~: algorithme, valeur, heure du calcul, témoin.
    \item Conditionnement~: enveloppe, sac antistatique, scellé, photo.
    \item Pour chaque transfert~: remettant, destinataire, date, heure,
          motif, état du scellé, signatures.
    \item Pour chaque examen~: personne, date, opération, copie utilisée.
    \item Conservation~: lieu, conditions (local fermé, armoire forte).
    \item Sort final~: restitution, destruction, archivage, avec
          autorisation.
\end{enumerate}
\end{boxprocedure}

\subsection{Le scellé}

Un scellé est un conditionnement qui ne s'ouvre pas sans laisser de marque~:
sac numéroté à fermeture inviolable, enveloppe signée sur les rabats et
recouverte d'un adhésif transparent. Le numéro du scellé figure dans le
formulaire. À chaque ouverture, on note qui l'a ouvert, pourquoi, et le
numéro du nouveau scellé. Dans la procédure pénale, la mise sous scellé
relève de l'officier de police judiciaire, comme le rappelle l'article~53 de
la loi n°~2010/012.

\begin{boxjuris}
\textbf{\loicam{n°~2010/012}, article~53}~\cite{loi2010012}~:
\textit{« (1) Les perquisitions en matière de cybercriminalité sont
susceptibles de porter sur des données qui peuvent être des supports
physiques ou des copies réalisées en présence des personnes qui assistent à
la perquisition. (2) Lorsqu'une copie des données saisies a été faite,
celle-ci peut être détruite sur instruction du Procureur de la République
pour des raisons de sécurité. (3) Sur accord du Procureur de la République,
seuls seront gardés sous scellé par l'Officier de Police Judiciaire, les
objets, documents et données utilisées à la manifestation de la vérité.
(4) Les personnes présentes lors de la perquisition peuvent être
réquisitionnées de fournir les renseignements sur les objets, documents et
données saisis. »}

\medskip
\textbf{\loicam{n°~2010/012}, article~54}~\cite{loi2010012}~:
\textit{« Les perquisitions et les saisies sont effectuées conformément aux
dispositions du Code de Procédure Pénale en tenant compte du dépérissement
des preuves. »}

\medskip
\textbf{Portée} (analyse). L'article~53 admet que la saisie porte sur une
copie des données, réalisée devant les personnes présentes~: c'est la base
légale de l'acquisition sur place, sans emporter le serveur. L'article~54
fait du dépérissement des preuves un critère de la conduite des saisies, ce
qui justifie de capturer d'abord ce qui disparaît le plus vite.
\end{boxjuris}

% ============================================================================
\section{La Chaîne de NYANGA}
\label{sec:c02-nyanga}
% ============================================================================

\begin{boxscenario}{Opération NYANGA. Le parcours du scellé NC-01}
\begin{tableaucours}{L{2.6cm}L{2.6cm}L{3.0cm}L{3.8cm}}{%
    \tableauHeader{Date et heure} &
    \tableauHeader{Détenteur} &
    \tableauHeader{Lieu} &
    \tableauHeader{Opération}
}
14/03, 10~h~02 & Mme NDOUMBE & Salle serveurs & Capture mémoire de
SRV-AD01 \\
14/03, 10~h~41 & Mme NDOUMBE & Salle serveurs & Empreinte SHA-256,
contresignée par M.~KAMGA \\
14/03, 15~h~50 & M. KAMGA & Bureau du DSI & Enveloppe kraft, ruban
adhésif, sans numéro \\
15/03, 9~h~30 & Mme MBARGA & Bureau du DSI & Empreinte recalculée,
identique~; mise sous scellé NC-01 \\
15/03, 11~h~00 & Mme MBARGA & Salle de crise & Copie de travail, empreinte
vérifiée \\
16/03, 10~h~15 & Adjudant-chef NANA & Brigade de recherches & Remise sur
réquisition, scellé intact, procès-verbal \\
17/03, 14~h~00 & M. ABOUBAKAR & Laboratoire de l'expert & Ouverture du
scellé, copie, nouveau scellé NC-01-E1 \\
\end{tableaucours}

\medskip
\textbf{Constatations.} La pièce a été détenue par cinq personnes
successives. Entre le 14~mars à 15~h~50 et le 15~mars à 9~h~30, soit
17~h~40, elle est restée dans un bureau ouvert, sans scellé numéroté et sans
consignation. Les empreintes calculées le 14 à 10~h~41 et le 15 à 9~h~30
sont identiques.

\textbf{Analyse} (confiance élevée sur l'intégrité, moyenne sur la
continuité). Le contenu n'a pas changé~: l'égalité des empreintes l'établit.
La continuité, elle, présente un trou. Il n'affaiblit pas l'intégrité, mais
il affaiblit la réponse à une autre objection~: la substitution d'un disque
par un autre de même contenu est sans effet, mais la manipulation du disque
pendant la nuit, par exemple pour en copier le contenu, ne peut être exclue.
Le rapport doit mentionner ce trou, et non le taire.
\end{boxscenario}

\begin{boxlecture}{Les six scellés de l'affaire au 17 mars}
\begin{tableaucours}{L{1.4cm}L{4.6cm}ccc}{%
    \tableauHeader{Scellé} &
    \tableauHeader{Pièce} &
    \tableauHeader{Intégrité} &
    \tableauHeader{Continuité} &
    \tableauHeader{Provenance}
}
NC-01 & Capture mémoire de SRV-AD01 & \OUI{} & \NEK{} & \OUI{} \\
NC-02 & Journaux de la passerelle (export) & \OUI{} & \OUI{} & \OUI{} \\
NC-03 & Image disque de SRV-CPT01 & \OUI{} & \OUI{} & \OUI{} \\
NC-04 & Capture réseau de la sonde & \OUI{} & \OUI{} & \OUI{} \\
NC-05 & Journaux de l'interface de paiement & \OUI{} & \OUI{} & \NEK{} \\
NC-06 & Capture d'écran WhatsApp d'un technicien & \NEK{} & \NON{} & \NEK{} \\
\end{tableaucours}

\medskip
\textbf{Lecture.} Sur dix-huit cotes, treize \OUI{}, une \NON{} et quatre
\NEK{}. NC-05 a été envoyé par le prestataire de l'interface sous forme de
fichier joint à un courriel~: son intégrité depuis l'envoi est établie, mais
rien ne dit comment le prestataire l'a extrait. NC-06 ne sert qu'à orienter
l'enquête~; on ne fonde aucune conclusion sur elle. NC-01 reste utilisable,
avec la réserve sur la nuit du 14 au 15.
\end{boxlecture}

% ============================================================================
\section{Les Ruptures et leurs Remèdes}
\label{sec:c02-ruptures}
% ============================================================================

\begin{tableaucours}{L{3.6cm}L{4.2cm}L{4.4cm}}{%
    \tableauHeader{Rupture} &
    \tableauHeader{Conséquence} &
    \tableauHeader{Remède}
}
Empreinte calculée tard & Rien ne prouve l'état d'avant & L'écrire, et
dater l'acquisition par d'autres sources (journal de l'outil) \\
Trou dans la détention & Une manipulation ne peut être exclue & Le dire
dans le rapport, recueillir les témoignages, vérifier les journaux d'accès
au local \\
Original examiné & L'original a pu changer & Recalculer l'empreinte,
documenter l'opération, travailler ensuite sur copie \\
Pièce transmise par courriel & Provenance invérifiable avant l'envoi &
Demander l'extraction en présence d'un témoin ou par procès-verbal \\
Scellé ouvert sans trace & La continuité est rompue & Nouveau scellé,
procès-verbal d'ouverture, signatures \\
\end{tableaucours}
\vspace{-0.8em}
\begin{center}{\small\itshape Tableau~2.3. Ruptures de chaîne et remèdes}\end{center}

Une rupture ne rend pas la pièce inutilisable. Elle en réduit la force
probante, d'autant plus que la pièce est seule à soutenir une conclusion.
Quand plusieurs pièces indépendantes convergent, la faiblesse de l'une est
compensée par les autres~: à NYANGA, l'outil d'accès à distance vu dans la
capture mémoire (NC-01) est aussi visible dans la capture réseau (NC-04),
par ses connexions vers 203.0.113.47. La conclusion ne repose pas sur NC-01
seule.

% ============================================================================
\section{Ce que la Preuve demande à chaque Métier}
\label{sec:c02-metiers}
% ============================================================================

\begin{boxmetier}{Application OPJ/Expert Judiciaire}
\textbf{Recevoir un scellé}
\begin{enumerate}
    \item Vérifier le numéro et l'état du scellé avant toute ouverture.
    \item Ouvrir en présence d'un témoin, photographier, consigner.
    \item Recalculer l'empreinte et la comparer à celle du formulaire.
    \item Travailler sur une copie, et resceller l'original.
    \item Signaler dans le rapport toute rupture constatée dans la chaîne
          reçue.
\end{enumerate}
\end{boxmetier}

\begin{boxmetier}{Analyste en réponse à incident}
Le formulaire de chaîne de possession s'ouvre dès le recueil, même quand
personne ne parle encore de plainte. Le coût est de quelques minutes par
pièce~; une chaîne reconstituée après coup ne convainc personne.
\end{boxmetier}

\begin{boxmetier}{RSSI}
Prévoir, avant l'incident, le matériel de conditionnement (sacs numérotés,
étiquettes, disques neufs) et un lieu de conservation fermé. À NYANGA, il
n'y avait ni l'un ni l'autre, d'où l'enveloppe kraft sur le bureau du DSI.
\end{boxmetier}

\separateur

\begin{boxresume}
\textbf{Ce qu'il faut retenir de ce chapitre~:}
\begin{itemize}
    \item La recevabilité dit si une pièce entre au débat~; la force
          probante dit ce qu'elle pèse. L'enquêteur construit la seconde.
    \item L'empreinte SHA-256 prouve que le contenu n'a pas changé entre
          deux calculs~; elle ne prouve ni la complétude, ni la provenance,
          ni l'état d'avant le premier calcul.
    \item La chaîne de possession retrace chaque détenteur, chaque
          transfert, chaque opération~; un trou se dit, il ne se cache pas.
    \item Les articles~53 et~54 de la loi n°~2010/012 admettent la saisie
          par copie et imposent de tenir compte du dépérissement des
          preuves.
    \item À NYANGA, NC-01 a une intégrité établie et une continuité \NEK{}
          sur 17~h~40~; la conclusion qu'elle soutient est confirmée par
          NC-04.
\end{itemize}
\end{boxresume}

\begin{boxinfo}
\textbf{Points de contrôle.} Vous avez compris ce chapitre si vous savez~:
\begin{itemize}
    \item expliquer pourquoi MD5 seul ne suffit plus, et ce qui reste vrai
          à son sujet~;
    \item remplir un formulaire de chaîne de possession~;
    \item coter l'intégrité, la continuité et la provenance d'une pièce.
\end{itemize}
\end{boxinfo}

% ============================================================================
\section*{Exercices}
% ============================================================================

\begin{boxexo}[title={\ding{45}\quad Exercice~2.1. Ce que dit l'empreinte \niveaubase}]
Vrai ou faux~? Justifiez.
\begin{enumerate}[(a)]
    \item Deux empreintes SHA-256 identiques prouvent que deux fichiers sont
          identiques.
    \item Une empreinte identique prouve que l'image disque est complète.
    \item Une image hachée en MD5 seulement est irrecevable.
    \item Il suffit de calculer l'empreinte au moment du rapport.
\end{enumerate}
\end{boxexo}

\begin{boxexo}[title={\ding{45}\quad Exercice~2.2. Remplir le formulaire \niveaumoyen}]
Le 15~mars à 16~h~00, SANKOFA EXPERTISE acquiert l'image du disque de
SRV-CPT01 (scellé NC-03). Rédigez le formulaire de chaîne de possession de
cette pièce jusqu'à sa remise à l'OPJ le 16~mars, en inventant de façon
cohérente les éléments manquants (outil, version, empreintes, témoins).
\end{boxexo}

\begin{boxexo}[title={\ding{45}\quad Exercice~2.3. La nuit du 14 au 15 \niveaumoyen}]
\begin{enumerate}[(a)]
    \item Rédigez le paragraphe du rapport qui mentionne le trou de la
          chaîne de NC-01, en séparant constatation et analyse.
    \item Quelles vérifications complémentaires proposez-vous pour réduire
          l'incertitude sur cette nuit~?
    \item Le trou change-t-il une conclusion de l'enquête~? Justifiez à
          l'aide de NC-04.
\end{enumerate}
\end{boxexo}

\begin{boxexo}[title={\ding{45}\quad Exercice~2.4. NC-05 et le prestataire \niveauexpert}]
Le prestataire de l'interface de paiement a envoyé ses journaux en pièce
jointe d'un courriel. L'avocat de la défense soutiendra qu'ils ont pu être
modifiés chez le prestataire.
\begin{enumerate}[(a)]
    \item Quelles cotes donnez-vous à cette pièce, et pourquoi~?
    \item Comment l'OPJ aurait-il dû obtenir ces journaux~? Appuyez-vous
          sur les articles~53 et~54.
    \item Quelles pièces indépendantes permettent de corroborer les heures
          des virements~?
\end{enumerate}
\end{boxexo}

\bigskip

\begin{boxinfo}
\textbf{Pour aller plus loin.}
\begin{itemize}
    \item L'ISO/IEC~27037~\cite{iso27037_2012}, sur l'identification, le
          recueil et la préservation.
    \item Le guide de l'ACPO~\cite{acpo2012}, pour ses quatre principes de
          conduite.
\end{itemize}
\end{boxinfo}

\begin{boxinfo}
\textbf{Transition.} La chaîne de possession garde la pièce~; encore
faut-il l'avoir acquise sans l'altérer. Copier un disque, capturer une
mémoire, figer un journal~: c'est l'objet du Chapitre~\ref{chap:acquisition}.
\end{boxinfo}
